% Generated draft: requires mathematical and source review.
\documentclass[11pt]{article}
\usepackage[margin=25mm]{geometry}
\usepackage{fontspec}
\setmainfont{Latin Modern Roman}
\usepackage{amsmath,amssymb,mathtools,unicode-math}
\setmathfont{Latin Modern Math}
\usepackage{enumitem,needspace,xurl,hyperref}
\setlength{\parindent}{0pt}
\setlength{\parskip}{5pt}
\setlength{\emergencystretch}{3em}
\newcommand{\N}{\mathbb{N}}
\newcommand{\Z}{\mathbb{Z}}
\newcommand{\Q}{\mathbb{Q}}
\newcommand{\R}{\mathbb{R}}
\newcommand{\C}{\mathbb{C}}
\newcommand{\F}{\mathbb{F}}
\newcommand{\A}{\mathbb{A}}
\newcommand{\PP}{\mathbb{P}}
\newcommand{\E}{\mathbb{E}}
\newcommand{\eps}{\varepsilon}
\newcommand{\dd}{\,\mathrm{d}}
\begin{document}
\section*{UnsolvedMath ID: 30006489}
\textbf{Source problem code:} catalogue entry
\subsection{Definitions and mathematical statement}
For a set \(X\), a map \(r\colon X\times X\to X\times X\) satisfies the set-theoretic Yang--Baxter equation in braid form if
\[
 r_{12}r_{23}r_{12}=r_{23}r_{12}r_{23}\qquad\text{on }X^3,
\]
where \(r_{12}=r\times\operatorname{id}_X\), \(r_{23}=\operatorname{id}_X\times r\), and composition is read from right to left. The supplied record does not specify whether the intended solutions must be bijective, nondegenerate, finite, or subject to other restrictions.

For an explicit rack convention, a right rack is a set \(X\) with an operation \(\triangleright\) such that every right translation \(x\mapsto x\triangleright y\) is bijective and
\[
 (x\triangleright y)\triangleright z
 = (x\triangleright z)\triangleright(y\triangleright z)
 \qquad(x,y,z\in X).
\]
Its associated rack solution is
\[
 r_{\triangleright}(x,y)=(y,x\triangleright y).
\]
This convention makes the rack terminology explicit; the supplied record does not identify its own left/right convention.

\textbf{Requested classification.} Classify set-theoretic Yang--Baxter solutions up to the intended notion of Drinfeld twist. In particular, for two rack solutions \(r,s\colon X^2\to X^2\) on the same set, determine whether
\[
 r\sim_{\mathrm{Dr}}s\quad\Longrightarrow\quad
 r(x,y)=s(x,y)\quad\text{for every }x,y\in X.
\]
Here \(\sim_{\mathrm{Dr}}\) is a placeholder for the twist relation intended by the source. Equality in the conclusion means equality of maps on the fixed set, as requested by the supplied wording. A precise definition of the admissible twists, and of the equivalence relation used for the general classification, is missing from the supplied record. Thus this is a conditional formulation of the requested problem, not a complete definition of its classification target.

\textbf{Qualification from the original excerpt.} The excerpt concerns a braided group \((G,m,1,r)\) and states that a group Drinfeld twist \((F,\Phi,\Psi)\) gives
\[
 G^F=\bigl(G,m\circ F^{-1},1,F\circ r\circ F^{-1}\bigr).
\]
In this construction the multiplication, as well as the braiding, changes. The excerpt supplies only part of the twist conditions and does not specify how this group construction determines the twist relation for arbitrary set-theoretic solutions. In particular, arbitrary conjugacy by a bijection of \(X^2\) must not be substituted for the missing twist definition. The original excerpt ends mid-sentence and does not itself contain the classification question quoted in the cleaned statement.
\subsection{Short English statement}
The supplied problem asks for a classification of set-theoretic Yang--Baxter solutions under Drinfeld twists, and asks whether two twist-related rack solutions on a fixed set must be exactly equal. The admissible twists and the intended class of solutions are not fully specified, so these qualifications remain unresolved. The supplied registry labels the problem open; that status has not been independently verified.
\subsection{Sources}
\begin{itemize}
\item Local UnsolvedMath registry, supplied record with permanent ID \(30006489\), exact source problem code \texttt{catalogue entry}, title \emph{Drinfeld-Twist Classification of Set-Theoretic Solutions}, in the Oberwolfach Reports Open Problems collection. No proposer is named; the supplied proposed year is 2025.
\item Supplied report citation: \emph{Mini-Workshop: The Yang--Baxter Equation and Representations of Braid Groups (2026)}, Oberwolfach Reports, pp.~2633--2672, DOI \texttt{10.4171/owr/2025/49}. \url{https://doi.org/10.4171/owr/2025/49}. The literature-source label is \emph{Oberwolfach Report 49/2025}; the supplied excerpt bears page number 2649. The 2026 year in the citation differs from the 2025 report designation and proposed year.
\item Supplied literature source: \emph{YB-semitrusses and set-theoretic solutions}. \url{https://arxiv.org/abs/2109.04978}.
\item The supplied literature assessment was checked on \texttt{2026-08-22} and labels the classification question open, including for rack solutions. This is supplied dated curation, not an independently verified literature assessment. The background explicitly requests independent verification of time-sensitive status and cited claims.
\item The registry labels the statement \texttt{corrected\_verified} and reports that the official report was checked and the question isolated and its notation repaired. These are supplied verification claims; no source visit or independent verification is asserted here.
\item The original excerpt attributes braided actions to De Commer, reference [2], and the assertion that group reflections yield group Drinfeld twists to Ferri, reference [5]. It also attributes the set-theoretic reflection consequence to reference [2]. Full bibliographic entries for these numbered references are not supplied.
\end{itemize}
\end{document}
